\section{Разработка архитектуры веб-приложения на основе Clean Architecture}

Архитектура проекта построена на принципах Clean Architecture, что позволяет разделить ответственность между слоями, минимизировать связанность и обеспечить расширяемость. В системе выделены два независимых сервиса: основной модуль задач (PostgreSQL) и модуль расписания (MongoDB). Каждый сервис имеет собственный слой API и внутренние компоненты.

\subsection{Слои приложения и их назначение}

Основное приложение (\texttt{app/}) разделено на следующие логические уровни:
\begin{enumerate}
    \item \textbf{Domain (сущности)} --- ORM-модели \texttt{User} и \texttt{Task}, определяющие структуру данных.
    \item \textbf{Application (бизнес-логика)} --- сервисы, реализующие правила обработки данных (\texttt{app/services/task\_service.py}).
    \item \textbf{Infrastructure (инфраструктура)} --- слой БД и конфигурации (\texttt{app/core}, \texttt{app/repositories}, \texttt{alembic}).
    \item \textbf{Presentation (представление)} --- FastAPI-роутеры и Jinja2-шаблоны (\texttt{app/routers}, \texttt{app/templates}).
\end{enumerate}

\subsection{Основные модули и зависимости}

Связи между слоями организованы по принципу зависимости внутрь:
\begin{itemize}
    \item \texttt{models/} описывает структуры БД и не зависит от API;
    \item \texttt{schemas/} описывает входные и выходные DTO для API и веб-форм;
    \item \texttt{repositories/} реализует CRUD-операции на уровне БД;
    \item \texttt{services/} агрегирует бизнес-логику и взаимодействует с репозиториями;
    \item \texttt{routers/} принимает HTTP-запросы и передает их в сервисы;
    \item \texttt{templates/} формируют HTML-страницы (SSR).
\end{itemize}

\subsection{Описание маршрута обработки запроса}

Типовой поток для API задач:
\begin{enumerate}
    \item Пользователь выполняет запрос к эндпоинту (например, \texttt{POST /tasks}).
    \item FastAPI роутер валидирует данные через Pydantic-схему \texttt{TaskCreate}.
    \item Сервис \texttt{task\_service.create\_task} формирует доменную сущность и вызывает репозиторий.
    \item Репозиторий \texttt{task\_repository.create} выполняет INSERT в PostgreSQL.
    \item Результат возвращается обратно в роутер и сериализуется в JSON.
\end{enumerate}

Для веб-интерфейса поток аналогичен, но результатом становится HTML-страница, отрендеренная Jinja2.

\subsection{Архитектура интеграции с сервисом расписания}

Сервис \texttt{rtu-mirea-schedule} развёрнут отдельно и предоставляет REST API (\texttt{/api/schedule/...}). Основной модуль использует эти данные на клиентской стороне (JS) для автодополнения предметов и отображения расписания. Это позволяет избежать прямой зависимости основного сервиса от MongoDB и внешних парсеров.

\subsection{Структура папок проекта}

Ниже приведена структура каталогов и назначение основных папок проекта (сокращённо, без служебных файлов):

\begin{verbatim}
курсач бекенд/
├── app/                      # основной FastAPI-сервис
│   ├── core/                 # конфигурация, БД, безопасность
│   ├── models/               # ORM-модели SQLAlchemy
│   ├── schemas/              # Pydantic-схемы (DTO)
│   ├── repositories/         # доступ к данным
│   ├── services/             # бизнес-логика
│   ├── routers/              # API и web-роуты
│   ├── templates/            # Jinja2-шаблоны
│   └── static/               # CSS и статические файлы
├── alembic/                  # миграции базы данных
│   └── versions/             # версии схемы БД
├── rtu-mirea-schedule/       # отдельный сервис расписания МИРЭА
│   ├── app/                  # FastAPI-сервис расписания
│   │   ├── api/              # API-эндпоинты
│   │   ├── core/             # конфигурация
│   │   ├── crud/             # операции MongoDB
│   │   ├── database/         # подключение к MongoDB
│   │   ├── models/           # Pydantic-модели расписания
│   │   └── schedule_parser/  # парсер расписания МИРЭА
│   └── docs/                 # документация сервиса
├── report/                   # LaTeX-отчёт по ГОСТ
│   ├── sections/             # разделы отчёта
│   └── bib/                  # библиография
├── docker-compose.yml        # PostgreSQL для локальной разработки
└── requirements.txt          # зависимости Python
\end{verbatim}
